\scanpage{46}
% archive foliation: 45
\uline{Cor} \deleted{$A$} Supposons $NA$ \unsure{nilp.t|nilpotent} (p. ex. $A$ noeth). Alors les conditions précédentes sur $u$ équivalent à :

(v) $u$ induit l'identité sur $A_{\mathrm{red}}$

Dans ces conditions, \deleted{\unsure{\exists}} si \deleted{\unsure{t}} $N(A)^{N+1}=0$, les conditions de la prop sont satisfaites pour $(u,N)$, i.e. $\forall$ opérateur diff de $A$ est d'ordre $\leq N$.

\uline{Cor.} Supposons $A$ une $K$-alg ($K$ une $\mathbb{Q}$-alg). Si $u=\exp\Theta$ ($\Theta$ dérivation nilpotente) alors $u$ est un $K$-autom ssi $\Theta$ est une $K$-dérivation.

\uline{Cor} Soit $I$ un idéal nilpotent de \deleted{$A$}, $K$-alg. comm. $A$ ($K$ $\mathbb{Q}$-alg. \unsure{commutative et unifère}\ldots), \inserted{$I^{N+1}=0$,} soit $L_K(A,I)$ l'ens des $K$-dérivations $\Theta$ \deleted{$A$} de $A$ qui induisent $0$ dans $A/I$ (i.e. $\Theta(A)\subset I$), $G_K(A,I)$ le grpe des $K$-autom. de $A$ \deleted{tels} qui induisent l'identité sur \deleted{\unsure{O}} $A/I$ (i.e. $(u-\operatorname{id})(A)\subset I$ i.e. $u(x)-x\in I$ $\forall x\in A$).

Alors a) $\forall\Theta\in L_K(A,I)$, on a $\Theta^{N+1}=0$

(plus généralement, $\forall\deleted{\unsure{k}}i\in\mathbb{N}$, $\Theta^i A\subset I^i$)

b) $\Theta\longmapsto\exp\Theta:\quad L_K(A,I)\simeq G_K(A,I)$
